% Capítulo "Inclusão" do relatório -- EDITÁVEL À MÃO (specs/2026-11-05_versao_final_site R5).
% Exportado uma vez por gera_latex.py --exportar-capitulos (2026-11-05); desde então o TEXTO é escrito aqui, não
% vem mais de relatorio/textos_curados.json (que é só do site). Os DADOS continuam gerados: as figuras
% (\figura{...}{../../visualizacoes/a4/...} e mapas/a4/) saem do analise.py; as tabelas no corpo
% (\input{gerado/tabelas_corpo/...}) e os rótulos do apêndice (\ref{tab:...}) saem do gera_latex.py, que confere
% a cada execução se tudo o que este arquivo cita existe.

\chapter{Inclusão}\label{cap:eixo-2}

\begin{achados}\begin{itemize}

\item Há aproximadamente 8 mil famílias com crianças com algum tipo de deficiência

\item Entre as cerca de 7,9 mil crianças com deficiência, os tipos mais registrados são transtorno mental (3,7 mil) e deficiência mental ou intelectual (3,3 mil); uma criança pode ter os dois

\item O BPC chega a cerca de 25\% das famílias com crianças com algum tipo de deficiência

\end{itemize}\end{achados}

Os dados de crianças com deficiência são oriundos exclusivamente do cadastro único. Dessa forma, não necessariamente refletem toda a primeira infância carioca, mas são relevantes para a parcela mais socialmente vulnerável. A predominância é de transtornos mentais e deficiência mental ou intelectual. Um percentual baixo (menos de 25\%) das famílias recebe Benefício de Prestação Continuada em função da deficiência. Como uma criança pode ter mais de um tipo de deficiência a soma do número de casos por tipo de deficiência ultrapassa o total de crianças.

\section{Crianças no CadÚnico com alguma deficiência}

\figura{mapa}{../../mapas/a4/mapa_percentual_cadunico_criancas_deficiencia_bairro_2026.pdf}{\% de crianças até 72 meses com deficiência no CadÚnico, por bairro}{Elaboração IPP com dados de \citeonline{mds_cadunico}. Limites de bairros: \citeonline{ipp_limites_bairros}. Sistema de referência SIRGAS 2000. Nota: Escala de cor limitada ao percentil 95 (5,1); o valor real está na tabela do apêndice.}{mapa:mapa-percentual-cadunico-criancas-deficiencia-bairro-2026}

A proporção de crianças com deficiência varia pouco entre os bairros: quase nove em cada dez ficam entre 3\% e 5\%, perto da média do município (3,9\%). Entre os bairros com valor próprio, os maiores percentuais estão em Pechincha (6,5\%), Itanhangá e Méier (6,0\% cada), Padre Miguel (5,3\%) e Caju (5,1\%). Pechincha e Méier têm só 25 crianças com deficiência cada; por isso pequenas diferenças no número mudam bastante o percentual. Os menores valores estão em Manguinhos (2,4\%), Cidade de Deus (2,5\%) e Acari (2,7\%). Em números absolutos, Campo Grande (446) e Santa Cruz (421) concentram mais crianças com deficiência. Para proteger os dados, 69 bairros com poucos casos aparecem somados a outros da mesma Região Administrativa e mostram a taxa do conjunto.

% Gerado por gera_latex.py -- não editar à mão.
{\small\setlength{\tabcolsep}{5pt}
\begin{longtable}{@{}>{\raggedright\arraybackslash}p{3.6cm}rrr@{}}
\caption{Crianças de 0 a 5 anos no CadÚnico com alguma deficiência}\label{tab:cadunico-deficiencia-criancas-0-a-5-2026}\\
\toprule \textbf{Indicador} & \textbf{Crianças} & \textbf{Base} & \textbf{\%} \\ \midrule \endfirsthead
\multicolumn{4}{@{}l}{\footnotesize\itshape (continuação)}\\ \toprule \textbf{Indicador} & \textbf{Crianças} & \textbf{Base} & \textbf{\%} \\ \midrule \endhead
\midrule \multicolumn{4}{r@{}}{\footnotesize\itshape (continua)}\\ \endfoot
\bottomrule \endlastfoot
Crianças até 72 meses no CadÚnico & 200.784 & -- & -- \\
Crianças com deficiência & 7.919 & 200.784 & 3,9 \\
\end{longtable}}
\vspace{-6pt}\fonte{Elaboração IPP com dados de \citeonline{mds_cadunico}.}


Das 200.784 crianças até 72 meses inscritas no CadÚnico em julho de 2026, 7.919 têm alguma deficiência registrada, o que corresponde a 3,9\% do total. O dado se refere apenas às crianças cadastradas e depende do que a família declarou no cadastro: não é uma estimativa da prevalência de deficiência entre todas as crianças da cidade.

\vertabelas{Tabela~\ref{tab:tabela-mapa-cadunico-deficiencia-bairro-2026}; ver Apêndice~\ref{ap:eixo-2}}

\section{Famílias no CadÚnico com criança com deficiência}

% Gerado por gera_latex.py -- não editar à mão.
{\small\setlength{\tabcolsep}{5pt}
\begin{longtable}{@{}>{\raggedright\arraybackslash}p{3.6cm}rrr@{}}
\caption{Famílias com crianças de 0 a 5 anos no CadÚnico: criança com deficiência e BPC por deficiência}\label{tab:cadunico-deficiencia-familias-0-a-5-2026}\\
\toprule \textbf{Indicador} & \textbf{Famílias} & \textbf{Base} & \textbf{\%} \\ \midrule \endfirsthead
\multicolumn{4}{@{}l}{\footnotesize\itshape (continuação)}\\ \toprule \textbf{Indicador} & \textbf{Famílias} & \textbf{Base} & \textbf{\%} \\ \midrule \endhead
\midrule \multicolumn{4}{r@{}}{\footnotesize\itshape (continua)}\\ \endfoot
\bottomrule \endlastfoot
Famílias com criança até 72 meses & 178.838 & -- & -- \\
Famílias com criança com deficiência & 7.810 & 178.838 & 4,4 \\
Delas, recebem BPC por deficiência & 1.530 & 6.335 & 24,2 \\
Delas, sem informação de BPC & 1.475 & -- & -- \\
\end{longtable}}
\vspace{-6pt}\fonte{Elaboração IPP com dados de \citeonline{mds_cadunico}.}


São 7.810 famílias com ao menos uma criança até 72 meses com deficiência, 4,4\% das 178.838 famílias com crianças nessa faixa no CadÚnico. Entre as 6.335 famílias com a informação preenchida, 1.530 (24,2\%) recebem o Benefício de Prestação Continuada (BPC) por deficiência; outras 1.475 não têm o dado. O BPC é registrado para a família, por isso o benefício pode se referir a qualquer membro, não necessariamente à criança.

\section{Crianças no CadÚnico por tipo de deficiência}

\figura{grafico}{../../visualizacoes/a4/cadunico_criancas_por_tipo_deficiencia.pdf}{CADÚNICO: Crianças até 72 meses com deficiência, por tipo}{Elaboração IPP com dados de \citeonline{mds_cadunico}.}{graf:cadunico-criancas-por-tipo-deficiencia}

Os tipos mais registrados são o transtorno mental (3.662 crianças, 46,2\% das crianças com deficiência) e a deficiência mental ou intelectual (3.332, 42,1\%), seguidos da deficiência física (2.035, 25,7\%) e da síndrome de Down (359, 4,5\%). As deficiências visuais e auditivas somam poucos casos, cada tipo abaixo de 2\%. Uma mesma criança pode ter mais de um tipo registrado, por isso os percentuais não somam 100\%.

\vertabelas{Tabela~\ref{tab:cadunico-tipos-deficiencia-0-a-5-2026}; ver Apêndice~\ref{ap:eixo-2}}

\section{Síntese do eixo}

\begin{sintese}{Inclusão}A amostra aponta que, mesmo num universo de vulnerabilidade social, há pouco alcance de benefícios sociais específicos. A co-ocorrência de diversas deficiências é também um desafio. Deficiências ligadas a surdez e baixa visão são menos frequentes, quando comparado a deficiências e transtornos mentais/intelectuais, indicando um ponto de atenção na disponibilização de serviços públicos.\end{sintese}

\eixo{3}{7}{users}
